\documentclass[10pt,a4paper]{amsart}
\usepackage[margin=19mm]{geometry}
\usepackage{amsmath,amssymb,mathtools}
\usepackage[scr=rsfs,cal=euler]{mathalfa}
\usepackage{xfrac}
\usepackage[unicode,hidelinks,bookmarks=false]{hyperref}
\newcommand{\ie}{i.e.\ }
\newcommand{\cf}{cf.\ }
\newcommand{\resp}{respectively\ }
\newcommand{\Subsection}{\subsection}
\providecommand{\nobreakdash}{\nobreak\hbox{-}}
\setlength{\emergencystretch}{2em}
\multlinegap0pt
\usepackage{amssymb}
\usepackage{mathtools}
\newtheorem{proposition}{Proposition}
\title{Shadowing property and transitivity of a set-valued map and its inverse limit}
\author{Yingcui Zhao, Lidong Wang$^\ast$}
\date{Source: arXiv:2607.17325v1 (CC BY 4.0)}
\begin{document}
\maketitle

\noindent\emph{Source and licence.} Shadowing property and transitivity of a set-valued map and its inverse limit. Version arXiv:2607.17325v1 (CC BY 4.0), distributed by arXiv under the Creative Commons Attribution 4.0 International licence, http://creativecommons.org/licenses/by/4.0/, as recorded in the arXiv licence metadata for this exact version and shown on the abstract page https://arxiv.org/abs/2607.17325v1. Full licence notice: \texttt{licenses/arxiv-2607.17325-CC-BY-4.0.txt}.\par\smallskip

\noindent\textit{Editorial scope.} Counted unit: the article's proposition prop1 (source0.tex lines 360--362). The article's complete original proof of it is delivered with it (source0.tex lines 363--390, one \texttt{proof} environment of 28 lines). No mathematics is written, repaired or re-derived: the statement and the proof are the article's own lines, in the article's own notation, and no retained line is substituted or re-flowed.  No further source-local context is needed: every cross-reference of every retained line resolves inside the two delivered spans.  Nothing was omitted from any retained span, so the package carries no omission record.  No retained line contains a citation, so the wrapper carries no bibliography and no bibliography entry.  No retained line contains a figure environment, an image-inclusion command or drawing code, so no image is delivered and the package declares no asset dependency.  Editorial formatting only, in addition: an AMS \texttt{amsart} preamble that restates the article's own declarations and macros used by the retained lines, namely the environments \texttt{proposition} on separate counters, together with the standard packages those lines need.  The retained environments print in this document as Proposition 1 in order of appearance, whereas the article prints the same items with the numbering of its own full document.  The complete native source of the article as distributed in the arXiv e-print for this exact version is bundled unchanged as \texttt{sources/205558/source0.tex}.
\begin{proposition}\label{prop1}
	Suppose that $X$ is a compact metric space with metric $d$ and $F:X\rightarrow 2^X$ is an upper semi-continuous set-valued map. Let $F(X)=X$. $F$ is transitive (resp. {weakly mixing, mixing, chain transitive, chain mixing}) if and only if $\underleftarrow{F}$ is transitive (resp. {weakly mixing, mixing, chain transitive, chain mixing}).
\end{proposition}

\begin{proof}
	{Based on $\underleftarrow{(\underleftarrow{F})}=F$, to prove this proposition, it is only necessary to prove necessity, that is, if $F$ is transitive (resp. {weakly mixing, mixing, chain transitive, chain mixing}) then $\underleftarrow{F}$ is transitive (resp. {weakly mixing, mixing, chain transitive, chain mixing}).}
	
	(1)Transitive: 
	
	{Let $U,V\subset X$ be two nonempty open sets. Since $F$ is transitive, there is $n\in\mathbb{N}$ such that $F^n(U)\bigcap V\neq\emptyset$. Then there is $x\in U$ such that one can find  $x_n\in F^n(x)\bigcap V$. Furthermore, there is $x\in\underleftarrow{F}^n(x_n)\bigcap U$ for some $x_n\in V$. That is, $\underleftarrow{F}^n(V)\bigcap U\neq\emptyset$. So, $\underleftarrow{F}$ is transitive.}
	
	(2)Weakly mixing:
	
	{Let $U,V\subset X$ be two nonempty open sets. Since $F$ is weakly mixing,  there is a strictly increasing subsequence $\{n_i\}$ of $\mathbb{N}$ such that for any $i\geq 1$ and any $0\leq j\leq i$, $F^{n_i+j}(U)\bigcap V\neq\emptyset$.
		Then,  $\underleftarrow{F}^{n_i+j}(V)\bigcap U\neq\emptyset$. So, $\underleftarrow{F}$ is weakly mixing.}
	
	(3)Mixing:
	Let $U,V\subset X$ be two nonempty open sets. Since $F$ is mixing, there is $M\in\mathbb{N}$ such that for any $n\geq M$, $F^{n}(U)\bigcap V\neq\emptyset$. Then, $\underleftarrow{F}^{n}(V)\bigcap U\neq\emptyset$. So, $\underleftarrow{F}$ is mixing.
	
	(4)Chain transitive:
	
	Let $x,y\in X$ and $\delta>0$. Since $\underleftarrow{F}$ is upper-continuous and $X$ is compact with the metric $d$, there is $\delta'>0$ such that for any $u\in X$, 
	\begin{equation}\label{ex1}
		\underleftarrow{F}(t)\subset B(\underleftarrow{F}(u),\delta), \forall t\in B(u,\delta').
	\end{equation}
	Since $F$ is chain transitive, there is a $\delta'$-chain $(x_i)^0_{i=n}$($n$ is a positive integer) of $F$ from $y$ to $x$ with $x_n=y$ and $x_0=x$. Then for any $0\leq i<n$, $d(F(x_{i+1}),x_{i})<\delta'$. That is, $F(x_{i+1})\bigcap B(x_i,\delta')\neq\emptyset$. Then $x_{i+1}\in\underleftarrow{F}(B(x_i,\delta'))$. By {Equation (\ref{ex1})}, $d(x_{i+1},\underleftarrow{F}(x_i))<\delta$. Then $(x_i)^n_{i=0}$ is a $\delta$-chain of $\underleftarrow{F}$ from $x$ to $y$. So, $\underleftarrow{F}$ is chain transitive.
	
	(5)Chain mixing:
	
	{Let $x,y\in X$ and $\delta>0$. Since $\underleftarrow{F}$ is upper-continuous and $X$ is compact with the metric $d$, there is $\delta'>0$ such that for any $u\in X$, $\underleftarrow{F}(t)\subset B(\underleftarrow{F}(u),\delta), \forall t\in B(u,\delta')$.
		Since $F$ is chain mixing, there is $N>0$ such that for any $n\geq N$ there exists a $\delta'$-chain $(x_i)^0_{i=n}$($n$ is a positive integer) of $F$ from $y$ to $x$. Then $(x_i)^n_{i=0}$ is a $\delta$-chain of $\underleftarrow{F}$ from $x$ to $y$. So, $\underleftarrow{F}$ is chain mixing.}
\end{proof}

\end{document}
